\documentclass{article}

\usepackage{amsmath,amssymb}
\usepackage{graphicx}

\title{Deep Learning Approaches for Natural Language Processing in Healthcare}

\author{Author Name}

\date{}

\begin{document}

\maketitle

\begin{abstract}
Artificial intelligence has transformed the healthcare industry in recent years. Deep learning models have shown remarkable performance in medical diagnosis tasks. Neural networks outperform traditional machine learning methods in image recognition tasks. This study proposes a transformer-based model for clinical text classification.
\end{abstract}

\textbf{Keywords:} Deep Learning, NLP, Healthcare, Transformers

\section{Introduction}
The adoption of electronic health records has increased significantly over the past decade. Studies have shown that clinical burnout is directly linked to documentation workload. Natural language processing can reduce administrative burden for physicians. Large language models have demonstrated human-level performance on medical licensing exams. The average hospital generates terabytes of unstructured text data annually. This paper presents a novel approach to automating clinical note classification.

\section{Methodology}
We collected data from three hospital systems and preprocessed using standard tokenization techniques.

\section{Results}
Our model achieved 94\% accuracy on the test set.

\section{Conclusion}
Deep learning shows strong promise for healthcare NLP applications. Further research is needed to address bias in training datasets.

\section*{References}
[1] Sample Ref.

\end{document}